\documentclass[11pt]{article}
\usepackage[margin=0.95in]{geometry}
\usepackage{amsmath}
\usepackage[hidelinks]{hyperref}
\pagestyle{empty}
\setlength{\parskip}{0.45em}
\setlength{\parindent}{0pt}

% DRAFT, 2026-11-14. Not sent. Items marked OWNER are for the author to
% fill in or decide before this letter is used.

\begin{document}

\begin{flushright}
Andrew H. Bond\\
Department of Computer Engineering\\
San Jos\'e State University\\
San Jos\'e, CA 95192 USA\\
andrew.bond@sjsu.edu\\[0.8em]
% OWNER: date of submission.
November 2026
\end{flushright}

To the Editor-in-Chief and Associate Editor\\
IEEE Transactions on Information Theory

Dear Editor:

Please consider the enclosed manuscript, ``Rate, Leakage, and Distortion
for Stationary Gaussian Sources: Water-Filling on a Tilted Spectrum,''
for publication in the IEEE Transactions on Information Theory. I am the
sole author. The review manuscript is in single-column format.

% OWNER: fill in the manuscript number of paper one.
This is the second of two manuscripts. The first, ``Rate, Leakage, and
Distortion for Gaussian Sources: Water-Filling on a Tilted Weight''
(reference [1] of the enclosed manuscript, manuscript number IT-26-XXXX),
is under review at this Transactions. It treats memoryless Gaussian
vector sources. The enclosed manuscript treats stationary Gaussian
vector processes, that is, sources with memory, and a second party
whose observation is filtered, delayed, or incomplete. It uses the
theorems of [1] by number and does not reprove them, so its results
depend on [1].
% OWNER: keep or drop the request in the next sentence.
For that reason the Editor may wish to assign it to the Associate
Editor who handles [1]. The theorem numbers of [1] that it cites are
those of the submitted version, and I will update them if [1] is
revised.

An encoder describes part of a stationary Gaussian vector process to a
decoder under a weighted quadratic distortion, while a second party
holds a noisy filtered observation of the process and never decodes. The
description has two costs per symbol, its rate and its conditional
leakage to the second party. With no second party this is reverse
water-filling for Gaussian processes (Kolmogorov [2], Berger [3]). I
found no rate--distortion--equivocation result for a Gaussian source
with memory and side information at the second party. The closest work
is that of Sankar, Rajagopalan, Mohajer, and Poor [14], [15], who
minimize an unconditional leakage for a scalar load in a circular block
model and note that the rate-optimal and leakage-optimal allocations
over frequency differ. The enclosed manuscript treats the two costs
jointly, with a leakage conditional on the second party's own
observation, for vector sources, and proves that the resulting region is
operational for the stationary source.

The main results are in three groups. First, the spectral region. For a
source of period $n$ the region splits exactly over the $n$ Fourier
frequencies under one common water level (Theorem~3), and as the period
grows it converges to a spectral region in which the tilted
water-filling equation of [1] holds at every frequency under one level
(Theorem~4). Second, explicit consequences. For a scalar process with a
scalar context the minimum leakage is an integral of the closed form of
[1], and the rate-optimal and leakage-optimal descriptions can occupy
disjoint frequency bands (Theorem~5 and Corollary~6). A delay or any
invertible filter at the second party changes nothing when it keeps its
whole observation (Proposition~7), and when it keeps less, the leakage
of a one-symbol description depends on what is kept through one
conditional variance (Proposition~8). Third, the stationary source,
which is the part that does not follow from [1] frequency by frequency.
For every continuous positive definite spectral density matrix and a
second party with a memoryless map and white noise, the finite-window
cost is at least the spectral cost at every window length and converges
to it, with a gap of at most a constant over the window length for
Lipschitz spectra (Theorem~15). The spectral region is the operational
region, with deterministic encoders, for that second party
(Theorem~16) and for one that observes the process through a filter in
noise of continuous positive definite spectrum, first for a filter of
finite length (Theorem~20) and then for every filter with a continuous
frequency response (Theorem~22).

The manuscript states what it does not prove. The sign and the bound
of the window gap are proved only for a memoryless map and white noise,
and the coding theorem for a filter gives no rate of convergence in the
blocklength.
The allocation over frequency in the scalar case is the allocation
among independent components of [1], with frequencies as the components,
and the manuscript says so.

Three points of disclosure. First, the dependence on [1] described
above. The memoryless coding theorem of [1], on which Theorems~16, 20, and~22
rest, is itself under review.
% OWNER: this sentence refers to manuscript IT-26-0996, as the cover
% letter of paper one did. Keep or drop.
Second, this Transactions declined a different manuscript of mine,
IT-26-0996, on 27 July 2026 without review, on scope and presentation.
That manuscript was a broad synthesis. The enclosed manuscript shares no
theorem with it. Third, no part of the enclosed manuscript has been
presented at a conference.

The manuscript's sources, check scripts, recorded outputs, and a
notebook that reruns every check are in a public repository, fixed for
this version by the tag \texttt{tit-stationary-1.0}
(\url{https://github.com/ahb-sjsu/geometric-observation/tree/tit-stationary-1.0/paper/tit-stationary}).
% OWNER: add the Zenodo DOI here if a deposit is made.
The statements are checked numerically on instances, independently of
the derivations, and the executed notebook records 53 passing checks and
no failure. The algebraic steps are machine-checked in a Lean~4 module
of twenty-one theorems. The appendix lists what is and is not
formalized. The limit theorems and the coding theorems are not
formalized, and no proof depends on the numerical checks.

One supplementary document is submitted as optional reading. It is a
short note that relates the results to a broader programme on
observer-dependent descriptions. The manuscript does not depend on it.

The use of AI tools in preparing the manuscript is disclosed in its
Acknowledgment, as IEEE policy requires. This manuscript is not under
review elsewhere, and I have no conflicts of interest to declare.

Sincerely,

Andrew H. Bond\\
Senior Member, IEEE

\end{document}
